% =============================================================================
% CAPÍTULO 2 — ESTADO DEL ARTE
% =============================================================================
\chapter{Estado del Arte}
\label{ch:estado_del_arte}

El presente capítulo revisa críticamente tres cuerpos de conocimiento relevantes para el desarrollo del sistema propuesto: los marcos normativos y procedimentales de la acreditación universitaria argentina; los avances en Sistemas Multiagente (MAS), inteligencia distribuida y coordinación cognitiva; y las aplicaciones recientes de Modelos de Lenguaje de Gran Escala (LLM), automatización documental y protocolos de interoperabilidad como el Model Context Protocol (MCP). El objetivo es identificar los aportes existentes, sus limitaciones y el vacío técnico-metodológico que justifica el desarrollo de este trabajo.

\section{Acreditación universitaria: normativa, debilidades y desafíos}
\label{sec:eda_acreditacion}

\subsection{Contexto histórico y evolución normativa}

La acreditación de carreras universitarias en Argentina surge como política de Estado a partir de la Ley de Educación Superior (LES) N°~24.521 de 1995, que en su artículo 43 establece la obligatoriedad de acreditar periódicamente aquellas carreras cuyo ejercicio pudiera comprometer el interés público. La misma ley crea la Comisión Nacional de Evaluación y Acreditación Universitaria (CONEAU) como organismo responsable de diseñar, coordinar y ejecutar los procesos de evaluación y acreditación en todo el sistema universitario \citep{marquina2022implicancias}. Desde sus inicios, la CONEAU desarrolló un modelo propio, flexible y diferenciado de los esquemas regionales y europeos, adaptado a la heterogeneidad institucional argentina \citep{greiner2023plan}.

En paralelo, el Consejo Federal de Decanos de Ingeniería (CONFEDI) impulsó durante más de tres décadas una transformación profunda del modelo formativo. La publicación del denominado ``Libro Rojo'' en 2018 estableció la educación basada en competencias como paradigma central para las carreras de ingeniería e informática, vinculado directamente a los estándares de acreditación. Este cambio exigió que cada asignatura articule resultados de aprendizaje esperados, metodologías de evaluación y nivel de aporte a las competencias de egreso, multiplicando la carga documental involucrada en los procesos de acreditación \citep{sanchez2022acreditacion, greiner2023plan}.

\subsection{Marco normativo vigente}

El procedimiento de acreditación de carreras en funcionamiento se rige actualmente por la 
Ordenanza CONEAU N°~62/17 \citep{ordenanza62} y la Ordenanza N°~75/23 \citep{ordenanza75}, 
las cuales actualizan y complementan el marco regulatorio vigente.

El proceso de acreditación se desarrolla a través de \textbf{cuatro etapas secuenciales}:

\begin{enumerate}
    \item \textbf{Formalización e inscripción en la convocatoria.}  
    La institución confirma su participación en el proceso de acreditación y registra la carrera 
    en la convocatoria correspondiente.

    \item \textbf{Autoevaluación institucional.}  
    La unidad académica realiza un análisis integral de la carrera y carga la información y 
    las evidencias correspondientes en el sistema \textit{CONEAU Global}.

    \item \textbf{Evaluación externa.}  
    Un \textbf{Comité de Pares Evaluadores} analiza la documentación presentada, pudiendo 
    realizar visitas institucionales y solicitar información complementaria.

    \item \textbf{Dictamen y resolución.}  
    El \textbf{plenario de la CONEAU} emite el dictamen final sobre la acreditación de la carrera.
\end{enumerate}

Cada una de estas etapas genera o requiere documentación específica, la cual debe organizarse 
en torno a \textbf{cuatro dimensiones de evaluación} establecidas por las resoluciones ministeriales:

\begin{table}[H]
  \centering
  \caption{Dimensiones de los estándares de acreditación (Resoluciones Ministeriales / Ordenanza CONEAU N°~62).}
  \label{tab:dimensiones_acreditacion}
  \begin{tabularx}{\textwidth}{lX}
    \toprule
    \textbf{Dimensión} & \textbf{Descripción} \\
    \midrule
    Contenidos Curriculares Básicos (Anexo I)    & Núcleo disciplinar mínimo y descriptores de conocimiento por terminal de ingeniería. \\
    Carga Horaria Mínima (Anexo II)              & Requisitos cuantitativos de horas por área y por actividad formativa. \\
    Criterios de Intensidad de la Formación Práctica (Anexo III) & Estándares cualitativos para actividades de práctica supervisada. \\
    Estándares para la Acreditación (Anexo IV)  & Criterios de evaluación por dimensión: contexto institucional, plan de estudios, cuerpo académico, infraestructura. \\
    \bottomrule
  \end{tabularx}
\end{table}

La carga de información en el sistema CONEAU Global constituye el punto de convergencia de toda la evidencia institucional, pero su completitud depende enteramente de la coordinación manual entre actores y departamentos dispersos \citep{greiner2023plan}.

\subsection{Complejidad organizacional del proceso}

Lo que los documentos normativos describen como un procedimiento lineal implica, en la práctica, una arquitectura organizacional de alta complejidad. \citet{greiner2023plan} documentan que en la Facultad de Ciencias Exactas de la UNNE la gestión del proceso requirió la creación de tres comisiones con composición heterogénea: la \textit{Comisión de Carrera} (responsable del plan de estudios), la \textit{Comisión Asesora de Cambio Curricular} (CACC, a cargo del diagnóstico de adecuación a los estándares) y la \textit{Comisión de Acreditación} (CA, responsable de la recolección sistemática de información para la carga en CONEAU Global). Estas comisiones operan en forma asincrónica, sin un sistema de integración que permita rastrear el estado, la versión ni la completitud de cada evidencia a lo largo del proceso.

Este patrón se repite en otras instituciones. \citet{sanabria2023procesos} analiza el proceso de acreditación de la carrera de Contador Público en tres universidades del noreste argentino y constata que los déficits reportados por la CONEAU en sus dictámenes se distribuyen en múltiples dimensiones de forma simultánea ---cuerpo docente, plan de estudios, infraestructura e investigación---, lo que refleja la dificultad de coordinar evidencias entre áreas que no comparten un repositorio común ni un lenguaje documental unificado.

\subsection{Desafíos identificados en la literatura}

A partir de la revisión de los trabajos citados, es posible sintetizar cuatro grandes desafíos que caracterizan el estado actual de la gestión de la acreditación universitaria en Argentina:

\begin{itemize}
    \item \textbf{Fragmentación y heterogeneidad documental:} los insumos para la autoevaluación provienen de múltiples fuentes (departamentos, áreas administrativas, sistemas institucionales locales) sin un repositorio integrador. En ausencia de herramientas de centralización, cada comisión gestiona sus evidencias de forma independiente, generando redundancias y versiones divergentes \citep{greiner2023plan, sanabria2023procesos}.

    \item \textbf{Falta de trazabilidad y control de versiones:} los documentos no suelen mantener un historial de cambios ni un registro claro de su vinculación con cada estándar evaluado. \citet{marquina2022implicancias}, en su encuesta a profesionales dedicados al aseguramiento de la calidad en universidades públicas argentinas, identifican que las dos tareas de mayor dedicación son el relevamiento de información y la redacción de informes ---actividades que se repiten en cada ciclo de acreditación sin acumulación de conocimiento estructurado entre períodos.

    \item \textbf{Altos costos temporales y cognitivos:} la lectura, clasificación y vinculación manual de evidencias demanda un esfuerzo sostenido de equipos especializados. \citet{marquina2022implicancias} reportan que los profesionales de aseguramiento de la calidad trabajan en equipos de aproximadamente seis personas con dedicación horaria fija, y que el principal obstáculo identificado es la ``falta de tiempo''. El 70\% califica sus tareas como ``novedosas o escasamente estructuradas'', lo que impide la estandarización y automati\-zación de los flujos de trabajo.

    \item \textbf{Baja reproducibilidad y cultura de cumplimiento formal:} distintos equipos clasifican y redactan informes de forma no homogénea, comprometiendo la consistencia entre ciclos y entre carreras. \citet{duarte2023literature}, en una revisión sistemática de sistemas de acreditación comparados, señala que los procesos pueden generar una cultura de cumplimiento superficial (\textit{compliance culture}) en la que las instituciones optimizan la presentación de indicadores en lugar de orientarse hacia la mejora genuina. \citet{casadiego2025acreditacion} agregan que la resistencia al cambio y la falta de formación continua en los equipos evaluadores condicionan la eficacia de los sistemas de acreditación.
\end{itemize}

\subsection{La brecha tecnológica}

Frente a la complejidad descrita, la literatura especializada coincide en señalar la ausencia de soporte tecnológico específico para el dominio de la acreditación como un factor estructural que limita la capacidad de las instituciones. \citet{sanabria2023procesos} concluye su investigación con una recomendación explícita: ``incorporar gestión tecnológica como factor clave para el logro de objetivos organizacionales de una manera eficaz y segura''. \citet{casadiego2025acreditacion} advierten que las tendencias emergentes ---digitalización, inteligencia artificial, internacionalización--- no están siendo incorporadas en los procesos de acreditación, que continúan operando con metodologías analógicas en un entorno que las supera.

En síntesis, la normativa vigente define con precisión \textbf{qué} debe evaluarse \citep{ordenanza62, ordenanza75}, pero no establece \textbf{cómo} gestionar tecnológicamente el flujo documental, la coordinación entre actores y el seguimiento del cumplimiento a lo largo del tiempo. Esta brecha entre exigencia normativa y capacidad operativa real define el problema que motiva el presente trabajo.

\section{Gestión documental inteligente}
\label{sec:eda_gestion_documental}

La gestión documental asistida por inteligencia artificial aborda un conjunto de problemas estructurales compartidos por numerosos dominios: grandes volúmenes de documentos heterogéneos, necesidad de clasificación semántica, trazabilidad de versiones y generación de reportes consistentes. Estos rasgos coinciden exactamente con los que caracterizan los procesos de acreditación universitaria descritos en la sección anterior, lo que convierte a este campo en un referente directo para el diseño de la plataforma propuesta.

\subsection{Pipeline tecnológico: de documentos no estructurados a información accionable}

El estado del arte en gestión documental inteligente converge hacia una arquitectura de \textbf{cinco etapas secuenciales}, organizadas en \textbf{tres capas funcionales}: (1)~\textbf{Ingesta}, que comprende la recepción de los documentos fuente y la extracción de texto mediante OCR o parsers de formato; (2)~\textbf{Procesamiento semántico}, donde modelos de lenguaje (NLP/LLM) identifican entidades, clasifican fragmentos y responden consultas; y (3)~\textbf{Recuperación}, que integra la indexación vectorial del contenido con la interfaz de consulta en lenguaje natural. La Figura~\ref{fig:pipeline_gdi} ilustra esta arquitectura con las tecnologías representativas de cada etapa.

\begin{figure}[H]
\centering
\begin{tikzpicture}[
    box/.style={
        rectangle, rounded corners=4pt,
        draw=Micolor1!80, fill=white,
        minimum width=2.0cm, minimum height=1.0cm,
        align=center, font=\small\bfseries},
    tech/.style={font=\scriptsize, text=gray!65, align=center},
    arr/.style={-{Stealth[length=5pt]}, thick, color=Micolor1!80},
    caplab/.style={font=\scriptsize\bfseries, align=center}]
  % Nodos principales
  \node[box, minimum width=2.8cm] (doc) {Documentos\\ no estructurados};
  \node[box, minimum width=2.8cm, right=0.55cm of doc]  (ocr) {OCR\\Extracción};
  \node[box, right=0.55cm of ocr]  (nlp) {NLP\,/\,LLM\\Procesamiento};
  \node[box, right=0.55cm of nlp]  (idx) {Indexación\\Vectorial};
  \node[box, right=0.55cm of idx]  (qry) {Consulta\\Semántica};
  % Flechas
  \draw[arr] (doc) -- (ocr);
  \draw[arr] (ocr) -- (nlp);
  \draw[arr] (nlp) -- (idx);
  \draw[arr] (idx) -- (qry);
  % Etiquetas tecnológicas (nombradas para los fondos de grupo)
  \node[tech, below=0.3cm of doc] (tdoc) {PDF / DOCX\\imágenes};
  \node[tech, below=0.3cm of ocr] (tocr) {Tesseract\\Google Vision};
  \node[tech, below=0.3cm of nlp] (tnlp) {Llama~3.1\\GPT / Mistral};
  \node[tech, below=0.3cm of idx] (tidx) {FAISS / Qdrant\\ChromaDB};
  \node[tech, below=0.3cm of qry] (tqry) {lenguaje\\natural};
  % Fondos de grupo que delimitan visualmente las tres capas
  \begin{pgfonlayer}{background}
    \fill[Micolor1!7, rounded corners=6pt]
      ([xshift=-22pt, yshift=-7pt]tdoc.south west) rectangle
      ([xshift=6pt,  yshift=19pt]ocr.north east);
    \fill[orange!10, rounded corners=6pt]
      ([xshift=-15pt, yshift=-7pt]tnlp.south west) rectangle
      ([xshift=6pt,  yshift=19pt]nlp.north east);
    \fill[teal!10, rounded corners=6pt]
      ([xshift=-6pt, yshift=-7pt]tidx.south west) rectangle
      ([xshift=6pt,  yshift=19pt]qry.north east);
  \end{pgfonlayer}
  % Etiquetas de capa sobre cada grupo
  \node[caplab, text=Micolor1!80]
    at ($(doc.north)!0.5!(ocr.north) + (0, 0.36cm)$)
    {\textit{Capa~1: Ingesta}};
  \node[caplab, text=orange!80]
    at ($(nlp.north) + (0, 0.36cm)$)
    {\textit{Capa~2: Procesamiento}};
  \node[caplab, text=teal!80]
    at ($(idx.north)!0.5!(qry.north) + (0, 0.36cm)$)
    {\textit{Capa~3: Recuperación}};
\end{tikzpicture}
\caption{Pipeline funcional de gestión documental inteligente: cinco etapas organizadas en tres capas. Elaboración propia basada en \citet{cascavita2025sistema} y \citet{torres2024desarrollo}.}
\label{fig:pipeline_gdi}
\end{figure}

\citet{cascavita2025sistema} implementan esta arquitectura completa integrando OCR (Tesseract, Google Vision API), modelos de lenguaje de código abierto (Llama~3.1, Phi-3) y bases vectoriales (FAISS, Qdrant). Sus resultados validan la viabilidad técnica del enfoque para organizaciones de tamaño mediano sin altos costos de licenciamiento, concluyendo que el sistema mejora la precisión en la consulta, reduce el tiempo de búsqueda y fortalece el cumplimiento normativo en la gestión documental. Los autores citan datos de McKinsey (2016) e IBM (2018) que cuantifican el alcance del problema: los empleados pueden destinar hasta el 20\% de su jornada laboral a buscar y gestionar información en documentos no estructurados, mientras que las organizaciones pueden perder hasta el 30\% de su eficiencia operativa por la gestión documental manual.

\citet{torres2024desarrollo} aplica este mismo pipeline al dominio de documentos legislativos del Parlamento de Canarias, un dominio regulatorio con marcadas similitudes al de la acreditación universitaria: documentos extensos, terminología normalizada, necesidad de clasificación temática y generación de resúmenes estructurados. El trabajo evalúa múltiples modelos de lenguaje (GPT, Mistral con fine-tuning), emplea ChromaDB como base vectorial para memoria semántica del sistema, y demuestra que la combinación de LLM y recuperación semántica permite responder consultas específicas sobre grandes corpus documentales de forma precisa y auditable.

\subsection{Aplicaciones en administración pública regulada}

El sector público representa un caso de estudio particularmente cercano al de la acreditación universitaria por sus requisitos de trazabilidad, transparencia y cumplimiento normativo. \citet{paun2025implementacion} desarrolla un marco metodológico y prototipo funcional para el Ayuntamiento de Valencia, enfocado en la extracción automatizada de datos desde solicitudes ciudadanas no estructuradas para su integración directa en el gestor de expedientes municipal. El trabajo identifica tres factores críticos de éxito para la implementación de IA en gestión documental pública:
\begin{itemize}
    \item trazabilidad completa del procesamiento,
    \item auditabilidad de cada decisión automatizada,
    \item integración con flujos de trabajo preexistentes sin disrupción.
\end{itemize}
Al mismo tiempo, documenta los principales desafíos que toda implementación real debe resolver:
\begin{itemize}
    \item \textbf{Técnicos:} precisión del OCR en documentos de baja calidad, errores de extracción NLP, interoperabilidad con sistemas heredados.
    \item \textbf{Organizativos:} resistencia al cambio, necesidad de formación del personal, supervisión manual en casos límite.
\end{itemize}

\citet{collado2024implementacion} estudia la aplicación de técnicas de Retrieval-Augmented Generation (RAG) sobre LLM para la extracción y generación de documentos en entidades públicas, demostrando que el patrón RAG mejora sustancialmente la precisión y pertinencia del contenido generado al anclar las respuestas del modelo en documentos reales del corpus institucional, en lugar de depender únicamente del conocimiento paramétrico del modelo. Este hallazgo es relevante para el presente trabajo, ya que las propuestas docentes y los estándares CONEAU constituyen exactamente el tipo de corpus normativo sobre el que un sistema RAG puede operar con alta fidelidad.

\subsection{Aplicaciones en monitoreo y seguimiento institucional}

\citet{tortosasistemas} presenta una revisión sistemática de la literatura sobre sistemas inteligentes para el seguimiento de proyectos ágiles en el contexto universitario argentino (UTN Resistencia), demostrando que los modelos de IA pueden extraer información relevante de los artefactos del proyecto, clasificar documentos de avance, identificar brechas respecto a los criterios de calidad definidos y sugerir acciones de mejora en tiempo real. El aspecto central de este trabajo para el presente contexto es la demostración de que la automatización del monitoreo ---es decir, el seguimiento del cumplimiento de criterios distribuidos en múltiples actores y dimensiones--- es técnicamente viable con las herramientas actuales de IA.

\subsection{Aplicaciones específicas en educación superior}

\citet{cruz2024diseno} describe el diseño y desarrollo de un sistema para la automatización de procesos institucionales en el Centro de Educación Continua de una universidad pública ecuatoriana. El trabajo implementa con metodología Scrum un sistema que automatiza los flujos de gestión de cursos, matrículas, certificaciones y reportes, con arquitectura Angular + Oracle y generación automatizada de diagramas de secuencia y modelos relacionales. Su contribución más relevante para el presente trabajo no es la solución técnica en sí misma, sino una advertencia metodológica que surge de sus conclusiones: la mayor dificultad en el desarrollo de sistemas de gestión universitaria no radica en la tecnología disponible, sino en la comprensión y modelado del marco normativo específico de cada institución. Esta observación adquiere una dimensión crítica cuando ese marco normativo es externo, como en el caso de la acreditación CONEAU, cuyos estándares evolucionan y deben reflejarse fielmente en el sistema para que este sea efectivo.

La Tabla~\ref{tab:gdi_comparativa} sintetiza los seis trabajos revisados, con sus dominios de aplicación, tecnologías centrales y el aporte específico que cada uno realiza al presente trabajo.

\begin{table}[htb]
\centering
\small
\setlength{\extrarowheight}{4pt}
\caption{Síntesis comparativa de trabajos en gestión documental inteligente.}
\label{tab:gdi_comparativa}
\arrayrulecolor{black}
\begin{tabularx}{\textwidth}{p{3.0cm}p{2.4cm}XX}
\toprule
\textbf{Trabajo} & \textbf{Dominio} & \textbf{Tecnología central} & \textbf{Aporte clave para este trabajo} \\
\midrule
\citet{cascavita2025sistema}      & Orgs.\ medianas       & OCR + Llama~3.1 + Qdrant/FAISS & Pipeline OCR-LLM-vectorial sin costo de licenciamiento \\
\arrayrulecolor{gray!35}\hline\arrayrulecolor{black}
\citet{torres2024desarrollo}      & Doc.\ legislativos    & GPT/Mistral + ChromaDB         & LLM+RAG sobre corpus regulatorios extensos \\
\arrayrulecolor{gray!35}\hline\arrayrulecolor{black}
\citet{paun2025implementacion}    & Adm.\ pública (ES)   & NLP + integración GDE          & Trazabilidad y auditabilidad como factores críticos \\
\arrayrulecolor{gray!35}\hline\arrayrulecolor{black}
\citet{collado2024implementacion} & Entidades públicas    & RAG sobre LLM                  & RAG mejora fidelidad sobre corpus normativos \\
\arrayrulecolor{gray!35}\hline\arrayrulecolor{black}
\citet{cruz2024diseno}            & Univ.\ (Ecuador)      & Angular + Oracle + Scrum       & Dificultad central: marco normativo institucional \\
\arrayrulecolor{gray!35}\hline\arrayrulecolor{black}
\citet{tortosasistemas}           & UTN Argentina         & IA para monitoreo ágil         & Automatización viable del seguimiento de criterios \\
\bottomrule
\end{tabularx}
\end{table}

\subsection{Brecha con el dominio CONEAU}

La revisión de los trabajos anteriores permite identificar un patrón consistente: los sistemas de gestión documental inteligente han demostrado eficacia técnica en dominios regulados (administración pública, parlamentos, universidades), pero ninguno de ellos aborda el ciclo de vida completo de la acreditación universitaria tal como lo define la normativa CONEAU. En particular, ninguna solución contempla de forma integrada:
\begin{itemize}
    \item el ciclo de vida de las propuestas de asignatura vinculado a los resultados de aprendizaje por competencias del Libro Rojo del CONFEDI,
    \item el control editorial diferenciado por roles institucionales (docente, director de carrera, secretaría académica),
    \item la trazabilidad de versiones de documentos como evidencia auditable ante el Comité de Pares evaluadores,
    \item la asistencia inteligente orientada específicamente a las dimensiones de los estándares de acreditación CONEAU.
\end{itemize}
Esta brecha técnica y de dominio es la que justifica el desarrollo propuesto en el presente trabajo.

\section{Sistemas Multiagente}
\label{sec:eda_mas}

Los \textbf{Sistemas Multiagente (MAS)} constituyen uno de los paradigmas centrales dentro del campo de la Inteligencia Artificial Distribuida. Un MAS está conformado por un conjunto de agentes autónomos que interactúan entre sí en un entorno compartido, cooperando para alcanzar objetivos colectivos mediante la coordinación de sus acciones individuales \citep{perez2022breve}. A diferencia de los sistemas monolíticos, los MAS permiten descomponer problemas complejos en unidades funcionales especializadas, distribuir la carga computacional y adaptar el comportamiento global ante cambios en el entorno.

\subsection{Fundamentos: tipos y propiedades de los agentes}

\citet{perez2022breve} sistematiza la tipología de agentes según sus capacidades de razonamiento en tres categorías fundamentales:

\begin{itemize}
    \item \textbf{Agentes reactivos}: responden directamente a estímulos del entorno sin mantener un modelo interno del mundo. Su fortaleza reside en la velocidad de respuesta y la robustez ante perturbaciones locales.
    \item \textbf{Agentes deliberativos}: mantienen una representación interna del entorno y razonan sobre ella para planificar acciones. El modelo más difundido es el \textbf{BDI} (\textit{Beliefs, Desires, Intentions}): el agente actualiza sus \textit{creencias} sobre el estado del mundo, las contrasta con sus \textit{deseos} y formula las \textit{intenciones} que guían su conducta.
    \item \textbf{Agentes híbridos}: combinan una capa reactiva para respuestas inmediatas con una capa deliberativa para planificación a largo plazo, equilibrando eficiencia y capacidad de razonamiento.
\end{itemize}

La Figura~\ref{fig:tipos_agentes} ilustra la arquitectura interna de cada tipo de agente, mostrando cómo varía la cadena \textbf{percepción--decisión--acción} según el nivel de razonamiento involucrado.

\tikzset{
  agbox/.style={rectangle, rounded corners=3pt,
                draw=Micolor1!75, fill=white,
                minimum width=1.6cm, minimum height=0.68cm,
                align=center, font=\tiny\bfseries},
  agenv/.style={rectangle, rounded corners=3pt,
                draw=gray!50, fill=gray!9,
                minimum width=1.4cm, minimum height=0.62cm,
                align=center, font=\tiny\itshape},
  agarr/.style={-{Stealth[length=3.5pt]}, semithick, color=Micolor1!75},
  agdarr/.style={-{Stealth[length=3.5pt]}, dashed, semithick, color=gray!50},
}

\begin{figure}[H]
\centering

%-------- (a) Reactivo --------
\begin{tikzpicture}
  \node[agenv]                      (env) {Entorno};
  \node[agbox, minimum width=2.8cm, below=0.5cm of env] (mod) {M\'{o}dulo reactivo};
  \node[agenv, below=0.5cm of mod]  (act) {Acci\'{o}n};
  \draw[agarr] (env) -- node[right,font=\small\itshape,text=gray!55]{percibe} (mod);
  \draw[agarr] (mod) -- node[right,font=\small\itshape,text=gray!55]{act\'{u}a} (act);
  \draw[agdarr] (act.west) -- ++(-0.8,0) |- (env.west);
  \begin{pgfonlayer}{background}
    \fill[Micolor1!6, rounded corners=5pt]
      ($(mod.west |- act.south) + (-0.65,-0.18)$) rectangle
      ($(mod.east |- env.north) + (0.65,0.18)$);
  \end{pgfonlayer}
\end{tikzpicture}\\[5pt]
{\small\textbf{(a) Agente reactivo} \,---\, \textit{sin estado interno: est\'{i}mulo $\to$ respuesta directa}}

\bigskip
%-------- (b) Deliberativo BDI --------
\begin{tikzpicture}
  \node[agenv, minimum width=2.8cm]                    (perc) {Percepci\'{o}n del entorno};
  \node[agbox, minimum width=2.8cm, below=0.5cm of perc] (bel) {Creencias \textit{(Beliefs)}};
  \node[agbox, minimum width=2.8cm, below=0.5cm of bel]  (des) {Deseos \textit{(Desires)}};
  \node[agbox, minimum width=2.8cm, below=0.5cm of des]  (int) {Intenciones \textit{(Intentions)}};
  \node[agenv, minimum width=2.8cm, below=0.5cm of int]  (act) {Acci\'{o}n};
  \draw[agarr] (perc) -- node[right,font=\small\itshape,text=gray!55]{actualiza} (bel);
  \draw[agarr] (bel)  -- node[right,font=\small\itshape,text=gray!55]{genera}    (des);
  \draw[agarr] (des)  -- node[right,font=\small\itshape,text=gray!55]{selecciona}(int);
  \draw[agarr] (int)  -- node[right,font=\small\itshape,text=gray!55]{ejecuta}   (act);
  \draw[agdarr] (act.east) -- ++(1.1,0) |- (perc.east)
    node[pos=0.18,right,font=\small\itshape,text=gray!55]{retroalim.};
  \begin{pgfonlayer}{background}
    \fill[orange!8, rounded corners=5pt]
      ($(bel.west |- act.south) + (-0.18,-0.18)$) rectangle
      ($(bel.east |- perc.north) + (0.18,0.18)$);
  \end{pgfonlayer}
\end{tikzpicture}\\[5pt]
{\small\textbf{(b) Agente deliberativo (BDI)} \,---\, \textit{razona sobre un modelo interno del mundo}}

\bigskip
%-------- (c) Híbrido --------
\begin{tikzpicture}
  \node[agenv, minimum width=2.8cm]            (env)   at (0,    0)     {Entorno};
  \node[agbox, minimum width=2.8cm, fill=orange!10] (delib) at (-1.7, -1.1)  {Capa deliberativa};
  \node[agbox, minimum width=2.8cm, fill=teal!8]    (react) at ( 1.7, -1.1)  {Capa reactiva};
  \node[agenv, minimum width=2.8cm]            (act)   at (0,   -2.2)   {Acci\'{o}n};
  \draw[agarr] (env.south west) -- node[left, font=\small\itshape,text=gray!55]{planifica} (delib.north);
  \draw[agarr] (env.south east) -- node[right,font=\small\itshape,text=gray!55]{responde}  (react.north);
  \draw[agarr] (delib.south) -- (act.north west);
  \draw[agarr] (react.south) -- (act.north east);
  \draw[agdarr] (act.east) -- ++(0.8,0) |- (env.east);
  \begin{pgfonlayer}{background}
    \fill[teal!7, rounded corners=5pt]
      ($(delib.west |- act.south) + (-0.18,-0.18)$) rectangle
      ($(react.east |- env.north) + (0.18,0.18)$);
  \end{pgfonlayer}
\end{tikzpicture}\\[5pt]
{\small\textbf{(c) Agente h\'{i}brido} \,---\, \textit{combina rapidez reactiva con planificaci\'{o}n deliberativa}}

\caption{Arquitecturas de agente según tipo de razonamiento. Elaboración propia basada en \citet{perez2022breve}.}
\label{fig:tipos_agentes}
\end{figure}

Independientemente de su tipo, los agentes comparten tres propiedades estructurales que los distinguen de los componentes de software tradicionales  \citep{sakurada2021multi}:
\begin{itemize}
    \item \textbf{Autonomía:} los agentes operan sin intervención humana continua, tomando decisiones y ejecutando acciones de forma independiente.
    \item \textbf{Cooperación:} los agentes coordinan sus acciones con otros agentes para alcanzar objetivos compartidos, mediante protocolos de comunicación y negociación.
    \item \textbf{Proactividad:} los agentes toman la iniciativa y no se limitan a reaccionar ante estímulos externos, sino que pueden generar planes y anticipar necesidades.
\end{itemize}

\subsection{Metodologías de diseño de sistemas MAS}

El diseño de un MAS requiere metodologías especializadas que capturen la dinámica de múltiples roles, organizaciones y normas de interacción. La Tabla~\ref{tab:mas_metodologias} sintetiza las principales metodologías identificadas en la literatura, siguiendo la clasificación de \citet{perez2022breve}.

\begin{table}[htb]
\centering
\small
\setlength{\extrarowheight}{4pt}
\caption{Principales metodologías de diseño de Sistemas Multiagente. Fuente: \citet{perez2022breve}.}
\label{tab:mas_metodologias}
\begin{tabularx}{\textwidth}{p{2.2cm}p{2.8cm}XX}
\toprule
\textbf{Metodología} & \textbf{Autor(es)} & \textbf{Enfoque principal} & \textbf{Fases del proceso} \\
\midrule
GAIA      & Wooldridge \& Jennings (2000)     & Diseño organizacional: roles, grupos y protocolos de interacción          & Análisis de roles $\to$ Modelos de interacción $\to$ Diseño arquitectónico \\
\arrayrulecolor{gray!35}\hline
TROPOS    & Bresciani et al.\ (2004)          & Desarrollo orientado a metas y actores con soporte para requisitos        & Requisitos tempranos $\to$ Arquitectura $\to$ Diseño detallado \\
\hline
PROMETHEUS & Padgham \& Winikoff               & Construcción de agentes inteligentes desde especificación formal          & Especificación $\to$ Diseño arquitectónico $\to$ Diseño detallado \\
\hline
PASSI     & Chella et al.\ (2006)             & Metodología paso a paso con notación UML; integra ingeniería de software  & Descripción del dominio $\to$ Identificación de agentes y roles $\to$ Tareas \\
\hline
RETSINA   & Graham et al.\ (2003)             & Arquitectura distribuida sin control centralizado; pares con ACL          & Comunidad de agentes $\to$ Infraestructura de comunicación $\to$ ACL \\
\arrayrulecolor{black}\bottomrule
\end{tabularx}
\end{table}

Un aspecto transversal a todas estas metodologías es la distinción entre el diseño del \emph{agente individual} y el de la \emph{organización}, que especifica roles, grupos, normas y estructuras sociales (jerarquías, coaliciones, federaciones) \citep{perez2022breve}. Esta separación resulta pertinente para el presente trabajo: los actores institucionales ---docentes, directores de carrera y secretaría académica--- se mapean naturalmente a roles diferenciados dentro de una organización MAS, con normativas explícitas que regulan sus responsabilidades y permisos.

\subsection{Coordinación y comunicación entre agentes}

La comunicación entre agentes es una propiedad constitutiva de los MAS y opera en cuatro dimensiones que deben contemplarse en el diseño \citep{perez2022breve}:

\begin{itemize}
    \item \textbf{Aspectos sociales}: conjunto de roles, grupos y relaciones entre ellos; las estructuras posibles incluyen jerarquías, coaliciones, equipos, mercados y organizaciones matriciales.
    \item \textbf{Aspectos comunicativos}: protocolos de intercambio de mensajes; los lenguajes de comunicación de agentes (ACL, FIPA) formalizan el contenido y la intención de cada mensaje.
    \item \textbf{Aspectos interactivos}: patrones de coordinación entre agentes (sincronización, negociación, votación).
    \item \textbf{Aspectos normativos}: obligaciones, prohibiciones y permisos que regulan el comportamiento; los agentes externos que deseen incorporarse a una organización deben acatar sus normas.
\end{itemize}

\citet{sakurada2021multi} extiende este marco al dominio industrial mostrando que los MAS permiten construir arquitecturas descentralizadas sobre el concepto de \textit{Asset Administration Shell} (AAS): una representación estandarizada de un activo que encapsula sus metadatos, versiones y relaciones en submodelos accesibles por agentes. Este concepto tiene un paralelo directo en el contexto de acreditación: cada propuesta de asignatura puede modelarse como un activo gestionado cuyo ciclo de vida ---creación, revisión, aprobación--- es coordinado por agentes con roles diferenciados.

\subsection{MAS con LLM: agentes cognitivos distribuidos}

La integración de Modelos de Lenguaje de Gran Escala en arquitecturas MAS representa el avance más reciente del paradigma. \citet{agashe2023llm} propone el marco \textbf{LLM-Co} (\textit{LLM-Coordination}), que define cómo un agente basado en LLM puede participar en tareas cooperativas sin protocolos de comunicación explícitos, mediante la capacidad de \textit{Theory of Mind}: la inferencia del estado mental del agente colaborador a partir del contexto compartido. La Figura~\ref{fig:llm_co_framework} ilustra la arquitectura interna del agente LLM-Co, organizada en tres capas funcionales.

\begin{figure}[H]
\centering
\begin{tikzpicture}[
    box/.style={
        rectangle, rounded corners=4pt,
        draw=Micolor1!80, fill=white,
        minimum width=2.2cm, minimum height=1.0cm,
        align=center, font=\small\bfseries},
    tech/.style={font=\scriptsize, text=gray!65, align=center},
    arr/.style={-{Stealth[length=5pt]}, thick, color=Micolor1!80},
    caplab/.style={font=\scriptsize\bfseries, align=center}]
  % Nodos principales
  \node[box] (env)   {Descripción\\ del entorno};
  \node[box, right=0.5cm of env]   (state) {Representación\\ de estado};
  \node[box, right=0.5cm of state] (feas)  {Acciones\\ factibles};
  \node[box, right=0.5cm of feas]  (llm)   {LLM\\ razonamiento};
  \node[box, right=0.5cm of llm]   (mgr)   {Gestor\\ de acciones};
  % Flechas
  \draw[arr] (env)   -- (state);
  \draw[arr] (state) -- (feas);
  \draw[arr] (feas)  -- (llm);
  \draw[arr] (llm)   -- (mgr);
  % Etiquetas tecnológicas
  \node[tech, below=0.3cm of env]   (tenv)   {\textit{Contexto}\\ \textit{del entorno}};
  \node[tech, below=0.3cm of state] (tstate) {\textit{Objetos}\\ \textit{y variables}};
  \node[tech, below=0.3cm of feas]  (tfeas)  {\textit{Acciones}\\ \textit{válidas}};
  \node[tech, below=0.3cm of llm]   (tllm)   {\textit{GPT-4}\\ \textit{LLaMA 3}};
  \node[tech, below=0.3cm of mgr]   (tmgr)   {\textit{Ejecución}\\ \textit{de acción}};
  % Fondos de grupo
  \begin{pgfonlayer}{background}
    \fill[Micolor1!7, rounded corners=6pt]
      ([xshift=-12pt, yshift=-7pt]tenv.south west) rectangle
      ([xshift=6pt,  yshift=19pt]state.north east);
    \fill[orange!10, rounded corners=6pt]
      ([xshift=-15pt, yshift=-7pt]tfeas.south west) rectangle
      ([xshift=6pt,  yshift=19pt]llm.north east);
    \fill[teal!10, rounded corners=6pt]
      ([xshift=-15pt, yshift=-7pt]tmgr.south west) rectangle
      ([xshift=6pt,  yshift=19pt]mgr.north east);
  \end{pgfonlayer}
  % Etiquetas de capa
  \node[caplab, text=Micolor1!80]
    at ($(env.north)!0.5!(state.north) + (0, 0.36cm)$)
    {\textit{Capa~1: Percepci\'{o}n}};
  \node[caplab, text=orange!80]
    at ($(feas.north)!0.5!(llm.north) + (0, 0.36cm)$)
    {\textit{Capa~2: Razonamiento}};
  \node[caplab, text=teal!80]
    at ($(mgr.north) + (0, 0.36cm)$)
    {\textit{Capa~3: Ejecuci\'{o}n}};
\end{tikzpicture}
\caption{Arquitectura del marco LLM-Co para agentes cognitivos en MAS. Elaboración propia basada en \citet{agashe2023llm}.}
\label{fig:llm_co_framework}
\end{figure}

\citet{bertelsmann_harnessing_multi_agent_2024} complementa esta visión desde el ecosistema de producción, describiendo el marco \textbf{AutoGen} \citep{wu2023autogen} como primera plataforma madura para la construcción de MAS generativos. Sus contribuciones más relevantes para este trabajo son: 
\begin{itemize}
    \item la asignación de diferentes configuraciones de LLM y temperatura según el rol y la tarea de cada agente;
    \item el concepto de \textit{group chat} como espacio de conocimiento compartido entre agentes;
    \item la descomposición automática de tareas complejas en subtareas asignadas a agentes especializados.
\end{itemize}

\citet{jimenez_romero_multi_agent_llm_2025} profundiza en los roles específicos que los LLM pueden cumplir dentro de un MAS: \textbf{razonamiento}, \textbf{mediación de decisiones}, \textbf{generación de explicaciones} y \textbf{clasificación semántica}.

\subsection{Brecha con el dominio CONEAU}

La literatura revisada confirma la madurez técnica del paradigma MAS y su extensión hacia sistemas basados en LLM. Sin embargo, \textbf{ninguno de los trabajos analizados aborda la aplicación de MAS al aseguramiento de la calidad universitaria bajo el marco normativo argentino}. Los roles institucionales involucrados en la acreditación ---docentes, directores de carrera, secretaría académica--- presentan características que los mapean naturalmente a agentes especializados: poseen conocimiento de dominio diferenciado, ejecutan tareas asincrónicas e interdependientes, y deben coordinar sus acciones respetando normas organizacionales explícitas. Esta correspondencia estructural entre el paradigma MAS y el proceso de acreditación CONEAU constituye uno de los fundamentos arquitectónicos del sistema propuesto.

\section{Modelos de Lenguaje de Gran Escala (LLM)}
\label{sec:eda_llm}

Los \textbf{Modelos de Lenguaje de Gran Escala (LLM)} representan uno de los avances más transformadores de la inteligencia artificial contemporánea. Se trata de modelos estadísticos de lenguaje preentrenados sobre enormes corpus textuales, capaces de comprender el contexto semántico, generar texto coherente, razonar sobre instrucciones complejas e interactuar con sistemas externos mediante lenguaje natural \citep{garcia2025desarrollo, farsane2025automatizacion}. Su irrupción desplazó paradigmas previos y habilitó una generación de sistemas que democratizan el acceso a capacidades analíticas antes reservadas a perfiles técnicos especializados.

\subsection{Arquitectura Transformer y fundamentos}

La base arquitectónica de todos los LLM modernos es el modelo \textbf{Transformer}, introducido por \citet{vaswani2017attention} en el artículo seminal \textit{Attention Is All You Need}. Su innovación central es el \textbf{mecanismo de atención}: en lugar de procesar las palabras secuencialmente, el modelo pondera simultáneamente la relevancia de cada token respecto a todos los demás, capturando dependencias de largo alcance con una eficiencia computacional sin precedentes \citep{vaswani2017attention, garcia2025desarrollo}.

Las Figuras~\ref{fig:transformer_arch} y \ref{fig:transformer_attention} reproducen las dos ilustraciones centrales del artículo original de \citet{vaswani2017attention}: la arquitectura \textbf{Encoder-Decoder} completa y el mecanismo de atención en sus variantes escalada y multi-cabeza.

%%────────────────────────────────────────────────────────────
%% FIGURA 1 – Arquitectura Transformer (Encoder + Decoder)
%%────────────────────────────────────────────────────────────
\begin{figure}[H]
\centering
\begin{tikzpicture}[
  blk/.style  ={rectangle, rounded corners=3pt,
                draw=Micolor1!70, fill=Micolor1!9,
                minimum width=3.0cm, minimum height=0.68cm,
                align=center, font=\scriptsize\bfseries},
  dblk/.style ={rectangle, rounded corners=3pt,
                draw=teal!70, fill=teal!9,
                minimum width=3.0cm, minimum height=0.68cm,
                align=center, font=\scriptsize\bfseries},
  nrm/.style  ={rectangle, rounded corners=2pt,
                draw=gray!40, fill=white,
                minimum width=3.0cm, minimum height=0.44cm,
                align=center, font=\tiny\itshape, text=gray!60},
  iob/.style  ={rectangle, rounded corners=3pt,
                draw=gray!50, fill=gray!8,
                minimum width=3.0cm, minimum height=0.55cm,
                align=center, font=\scriptsize},
  emb/.style  ={rectangle, rounded corners=3pt,
                draw=orange!60, fill=orange!7,
                minimum width=3.0cm, minimum height=0.55cm,
                align=center, font=\scriptsize},
  outb/.style  ={rectangle, rounded corners=3pt,
                draw=orange!60, fill=orange!7,
                minimum width=3.0cm, minimum height=0.55cm,
                align=center, font=\scriptsize},
  sof/.style  ={rectangle, rounded corners=3pt,
                draw=gray!45, fill=gray!6,
                minimum width=3.0cm, minimum height=0.55cm,
                align=center, font=\scriptsize},
  arr/.style  ={-{Stealth[length=3.5pt]}, semithick, color=gray!65},
  rarr/.style ={-{Stealth[length=3pt]}, semithick, color=gray!50, dashed},
  Narr/.style ={-{Stealth[length=3.5pt]}, semithick, color=teal!60, dashed},
]

%% ═══════════════ ENCODER (izquierda) ═══════════════
\node[iob]  (EI)  at (0, 0.00) {Inputs};
\node[emb]  (EE)  at (0, 0.90) {Input Embedding};
\node[iob]  (EP)  at (0, 1.75) {Positional\\Encoding};
\node[blk]  (EA)  at (0, 3.00) {Multi-Head\\Attention};
\node[nrm]  (EN1) at (0, 3.90) {Add \& Norm};
\node[blk]  (EF)  at (0, 4.90) {Feed Forward};
\node[nrm]  (EN2) at (0, 5.80) {Add \& Norm};

\draw[arr] (EI)  -- (EE);
\draw[arr] (EE)  -- (EP);
\draw[arr] (EP)  -- (EA.south);
\draw[arr] (EA)  -- (EN1);
\draw[arr] (EN1) -- (EF);
\draw[arr] (EF)  -- (EN2);

% Residual 1 encoder
\draw[rarr] (0, 2.60) -- (-2.1, 2.60) -- (-2.1, 3.90) -- (EN1.west);
\fill[gray!55] (0, 2.60) circle (1.4pt);
% Residual 2 encoder
\draw[rarr] (0, 4.50) -- (-2.1, 4.50) -- (-2.1, 5.80) -- (EN2.west);
\fill[gray!55] (0, 4.50) circle (1.4pt);

% Brace ×N encoder
\node[font=\tiny\bfseries, text=Micolor1!70] at (-2.5, 4.40) {$\times N$};

%% ═══════════════ DECODER (derecha) ═══════════════
\node[iob]  (DI)  at (6.5, 0.00) {Outputs \small(shifted right)};
\node[emb]  (DE)  at (6.5, 0.90) {Output Embedding};
\node[iob]  (DP)  at (6.5, 1.75) {Positional\\Encoding};
\node[dblk] (DA1) at (6.5, 3.00) {Masked\\Multi-Head Attention};
\node[nrm]  (DN1) at (6.5, 3.90) {Add \& Norm};
\node[dblk] (DA2) at (6.5, 5.00) {Multi-Head\\Attention};
\node[nrm]  (DN2) at (6.5, 5.90) {Add \& Norm};
\node[dblk] (DF)  at (6.5, 6.90) {Feed Forward};
\node[nrm]  (DN3) at (6.5, 7.80) {Add \& Norm};
\node[sof]  (DS)  at (6.5, 8.70) {Linear + Softmax};
\node[outb]  (DO)  at (6.5, 9.55) {Output\\Probabilities};

\draw[arr] (DI)  -- (DE);
\draw[arr] (DE)  -- (DP);
\draw[arr] (DP)  -- (DA1.south);
\draw[arr] (DA1) -- (DN1);
\draw[arr] (DN1) -- (DA2.south);
\draw[arr] (DA2) -- (DN2);
\draw[arr] (DN2) -- (DF);
\draw[arr] (DF)  -- (DN3);
\draw[arr] (DN3) -- (DS);
\draw[arr] (DS)  -- (DO);

% Residual 1 decoder
\draw[rarr] (6.5, 2.60) -- (8.9, 2.60) -- (8.9, 3.90) -- (DN1.east);
\fill[gray!55] (6.5, 2.60) circle (1.4pt);
% Residual 2 decoder
\draw[rarr] (6.5, 4.55) -- (8.9, 4.55) -- (8.9, 5.90) -- (DN2.east);
\fill[gray!55] (6.5, 4.55) circle (1.4pt);
% Residual 3 decoder
\draw[rarr] (6.5, 6.48) -- (8.9, 6.48) -- (8.9, 7.80) -- (DN3.east);
\fill[gray!55] (6.5, 6.48) circle (1.4pt);

% Brace ×N decoder
\node[font=\tiny\bfseries, text=teal!70] at (9.3, 5.40) {$\times N$};

%% ═══ Flecha Encoder → Decoder (cross-attention) ═══
\draw[Narr, line width=0.8pt]
  (EN2.north) -- (0, 6.60) -- (0, 7.30) -- (3.3, 7.30) --
  node[above, font=\scriptsize\itshape, text=teal!60]{K, V} (DA2.west);
% segunda línea K/V
\draw[Narr, line width=0.8pt]
  (0, 6.60) -- (3.3, 6.60) -- (3.3, 7.30);
\fill[teal!50] (0, 6.60) circle (1.4pt);

%% ═══ Fondos ═══
\begin{pgfonlayer}{background}
  \fill[Micolor1!5, rounded corners=7pt]
    (-2.35, -0.35) rectangle (2.35, 6.35);
  \fill[teal!5, rounded corners=7pt]
    (4.15, -0.35) rectangle (9.10, 9.90);
\end{pgfonlayer}

%% Etiquetas columnas
\node[font=\tiny\bfseries, text=Micolor1!75] at (0,   6.85) {Encoder};
\node[font=\tiny\bfseries, text=teal!75]     at (6.5, 10.15) {Decoder};

\end{tikzpicture}
\caption[Arquitectura Encoder-Decoder del modelo Transformer.]{Arquitectura Encoder-Decoder del modelo Transformer. Las flechas en L representan conexiones residuales; las flechas punteadas K,\,V indican la atención cruzada (\textit{cross-attention}) del decoder sobre las representaciones del encoder. \footnotesize\textit{Elaboración propia basada en} \protect\citeauthor{vaswani2017attention} (\protect\citeyear{vaswani2017attention}).}
\label{fig:transformer_arch}
\end{figure}

%%────────────────────────────────────────────────────────────
%% FIGURA 2 – Scaled Dot-Product + Multi-Head Attention
%%────────────────────────────────────────────────────────────
\begin{figure}[H]
\centering
\resizebox{0.65\textwidth}{!}{%
\begin{tikzpicture}[
  op/.style  ={circle, draw=gray!55, fill=white,
               minimum size=0.55cm, align=center, font=\scriptsize},
  blk/.style ={rectangle, rounded corners=3pt,
               draw=Micolor1!65, fill=Micolor1!8,
               minimum width=2.2cm, minimum height=0.60cm,
               align=center, font=\scriptsize\bfseries},
  linb/.style={rectangle, rounded corners=3pt,
               draw=teal!60, fill=teal!7,
               minimum width=2.2cm, minimum height=0.55cm,
               align=center, font=\scriptsize},
  iob/.style ={rectangle, rounded corners=3pt,
               draw=gray!45, fill=gray!7,
               minimum width=1.4cm, minimum height=0.50cm,
               align=center, font=\scriptsize},
  mhb/.style ={rectangle, rounded corners=3pt,
               draw=Micolor1!65, fill=Micolor1!8,
               minimum width=2.2cm, minimum height=0.60cm,
               align=center, font=\scriptsize\bfseries},
  arr/.style ={-{Stealth[length=3pt]}, semithick, color=gray!60},
  garr/.style={-{Stealth[length=3pt]}, semithick, color=Micolor1!60},
]

%% ═══ Columna A: Scaled Dot-Product Attention ═══
\node[iob]  (AQ) at (0.0, 0.0) {Q};
\node[iob]  (AK) at (1.1, 0.0) {K};
\node[iob]  (AV) at (2.2, 0.0) {V};

\node[blk]  (AMM) at (0.55, 1.05) {MatMul};
\node[blk]  (ASC) at (0.55, 2.00) {Scale ($\div\sqrt{d_k}$)};
\node[blk, draw=gray!45, fill=gray!7]
            (AMK) at (0.55, 2.95) {Mask (opt.)};
\node[blk]  (ASF) at (0.55, 3.90) {SoftMax};
\node[blk]  (AM2) at (1.10, 5.00) {MatMul};
\node[iob]  (AO)  at (1.10, 6.00) {Output};

\draw[arr] (AQ.north)  |- (AMM.south);
\draw[arr] (AK.north)  |- (AMM.south);
\draw[arr] (AMM) -- (ASC);
\draw[arr] (ASC) -- (AMK);
\draw[arr] (AMK) -- (ASF);
\draw[arr] (ASF.north) |- (AM2.south);
\draw[arr] (AV.north)  |- (AM2.south);
\draw[arr] (AM2) -- (AO);

% Fondo

\begin{pgfonlayer}{background}
  \fill[Micolor1!5, rounded corners=6pt]
    (-0.85, -0.35) rectangle (3.2, 6.35);
\end{pgfonlayer}
\node[font=\scriptsize\bfseries, text=Micolor1!70, align=center]
  at (1.10, 6.75) {Scaled Dot-Product\\Attention};

%% ═══ Columna B: Multi-Head Attention ═══
\node[iob]  (BQ) at (5.5, 0.0) {Q};
\node[iob]  (BK) at (6.6, 0.0) {K};
\node[iob]  (BV) at (7.7, 0.0) {V};

\node[linb] (BLQ) at (5.5, 1.1) {Linear};
\node[linb] (BLK) at (6.6, 1.1) {Linear};
\node[linb] (BLV) at (7.7, 1.1) {Linear};

% Bloque "Scaled Dot-Product ×h" como caja grande
\node[mhb, minimum width=3.5cm, minimum height=1.0cm]
            (BSP)  at (6.6, 2.65) {Scaled Dot-Product\\Attention $\times h$};
\node[linb, minimum width=3.5cm]
            (BCC)  at (6.6, 3.90) {Concat};
\node[linb, minimum width=3.5cm]
            (BLO)  at (6.6, 4.80) {Linear};
\node[iob]  (BO)   at (6.6, 5.85) {Output};

\draw[arr] (BQ.north) -- (BLQ);
\draw[arr] (BK.north) -- (BLK);
\draw[arr] (BV.north) -- (BLV);
\draw[arr] (BLQ.north) |- (BSP.south);
\draw[arr] (BLK.north) -- (BSP.south);
\draw[arr] (BLV.north) |- (BSP.south);
\draw[arr] (BSP) -- (BCC);
\draw[arr] (BCC) -- (BLO);
\draw[arr] (BLO) -- (BO);

\begin{pgfonlayer}{background}
  \fill[teal!5, rounded corners=6pt]
    (4.55, -0.35) rectangle (8.8, 6.35);
\end{pgfonlayer}
\node[font=\scriptsize\bfseries, text=teal!70, align=center]
  at (6.6, 6.75) {Multi-Head Attention};

%% Etiqueta h cabezas
\node[font=\tiny\itshape, text=gray!55, align=center]
  at (9.1, 2.65) {$h$\\cabezas\\paralelas};

\end{tikzpicture}
}% fin resizebox
\caption[Mecanismos de atención del Transformer.]{Mecanismos de atención del Transformer. \textbf{(izquierda)} \textit{Scaled Dot-Product Attention}: $\mathrm{Attention}(Q,K,V)=\mathrm{softmax}(QK^{\top}/\sqrt{d_k})\,V$. \textbf{(derecha)} \textit{Multi-Head Attention}: $h$ cabezas de atención en paralelo cuyos resultados se concatenan y proyectan linealmente. \footnotesize\textit{Elaboración propia basada en} \protect\citeauthor{vaswani2017attention} (\protect\citeyear{vaswani2017attention}).}
\label{fig:transformer_attention}
\end{figure}

Sobre esta base surgieron tres familias de modelos según su configuración arquitectónica:

\begin{itemize}
    \item \textbf{Encoder-only} (ej.\ BERT, 2018): procesan toda la secuencia de entrada de forma bidireccional; ideales para tareas de comprensión como clasificación de texto y extracción de entidades.
    \item \textbf{Decoder-only} (ej.\ GPT, LLaMA, Mistral): generan texto de izquierda a derecha prediciendo el siguiente token a partir del contexto acumulado; son la base de los asistentes conversacionales actuales.
    \item \textbf{Encoder-decoder} (ej.\ T5, BART): combinan ambas variantes; adecuados para traducción automática y resumen extractivo.
\end{itemize}

La Figura~\ref{fig:transformer_familias} ilustra de forma esquemática las diferencias arquitectónicas entre las tres familias, incluyendo la dirección del mecanismo de atención, los modelos representativos y sus aplicaciones principales.

\begin{figure}[H]
\centering
\begin{tikzpicture}[
  hdr/.style={rectangle, rounded corners=3pt,
              minimum width=3.4cm, minimum height=0.65cm,
              align=center, font=\small\bfseries, text=white},
  mblk/.style={rectangle, rounded corners=3pt,
               minimum width=3.0cm, minimum height=1.35cm,
               align=center, font=\scriptsize},
  hblk/.style={rectangle, rounded corners=3pt,
               minimum width=3.0cm, minimum height=0.65cm,
               align=center, font=\scriptsize},
  iobox/.style={rectangle, rounded corners=2pt,
                draw=gray!55, fill=gray!7,
                minimum width=3.0cm, minimum height=0.52cm,
                align=center, font=\scriptsize},
  parr/.style={-{Stealth[length=3.5pt]}, semithick, color=gray!55},
  xarr/.style={-{Stealth[length=3.5pt]}, semithick, color=orange!65, dashed},
  nbox/.style={align=center, font=\scriptsize, text=gray!65}
]

%% ─── Columna 1: Encoder-only (burdeos) ─────────────────────
\node[hdr, fill=Micolor1!80]                    (h1) at (1.7,  5.0) {Encoder-only};
\node[iobox]                                    (i1) at (1.7,  4.2) {\textit{Entrada}};
\node[mblk, draw=Micolor1!65, fill=Micolor1!8] (b1) at (1.7,  2.85)
    {\textbf{Encoder} $\times N$\\[4pt]
     Atenci\'{o}n\\bidireccional};
\node[iobox]                                    (o1) at (1.7,  1.5) {\textit{Representaci\'{o}n}};
\node[nbox]                                     (n1) at (1.7,  0.7)
    {BERT $\cdot$ RoBERTa\\[1pt]Clasificaci\'{o}n $\cdot$ NER};
\draw[parr] (i1.south) -- (b1.north);
\draw[parr] (b1.south) -- (o1.north);

%% ─── Columna 2: Decoder-only (teal) ────────────────────────
\node[hdr, fill=teal!65]                        (h2) at (6.1,  5.0) {Decoder-only};
\node[iobox]                                    (i2) at (6.1,  4.2) {\textit{Tokens previos}};
\node[mblk, draw=teal!65, fill=teal!8]          (b2) at (6.1,  2.85)
    {\textbf{Decoder} $\times N$\\[4pt]
     Atenci\'{o}n causal\\(izq.\ $\rightarrow$ der.)};
\node[iobox]                                    (o2) at (6.1,  1.5) {\textit{Siguiente token}};
\node[nbox]                                     (n2) at (6.1,  0.7)
    {GPT $\cdot$ LLaMA $\cdot$ Mistral\\[1pt]Generaci\'{o}n $\cdot$ Di\'{a}logo};
\draw[parr] (i2.south) -- (b2.north);
\draw[parr] (b2.south) -- (o2.north);

%% ─── Columna 3: Encoder-Decoder (orange) ───────────────────
\node[hdr, fill=orange!70]                      (h3) at (10.5, 5.0) {Encoder-Decoder};
\node[iobox]                                    (i3) at (10.5, 4.2) {\textit{Texto fuente}};
\node[hblk, draw=orange!65, fill=orange!8]      (e3) at (10.5, 3.3)
    {\textbf{Encoder} $\times N$ --- Bidir.};
\node[hblk, draw=orange!65, fill=orange!8]      (d3) at (10.5, 2.35)
    {\textbf{Decoder} $\times N$ --- Causal};
\node[iobox]                                    (o3) at (10.5, 1.5) {\textit{Texto generado}};
\node[nbox]                                     (n3) at (10.5, 0.7)
    {T5 $\cdot$ BART\\[1pt]Traducci\'{o}n $\cdot$ Resumen};
\draw[parr] (i3.south) -- (e3.north);
\draw[xarr] (e3.south) --
    node[right, font=\tiny\itshape, text=orange!65]{cross-attn} (d3.north);
\draw[parr] (d3.south) -- (o3.north);

%% ─── Fondos de grupo ────────────────────────────────────────
\begin{pgfonlayer}{background}
  \fill[Micolor1!6, rounded corners=7pt]
    ([xshift=-14pt, yshift=-9pt]n1.south -| h1.west) rectangle
    ([xshift= 14pt, yshift= 7pt]h1.north -| h1.east);
  \fill[teal!6, rounded corners=7pt]
    ([xshift=-14pt, yshift=-9pt]n2.south -| h2.west) rectangle
    ([xshift= 14pt, yshift= 7pt]h2.north -| h2.east);
  \fill[orange!6, rounded corners=7pt]
    ([xshift=-14pt, yshift=-9pt]n3.south -| h3.west) rectangle
    ([xshift= 14pt, yshift= 7pt]h3.north -| h3.east);
\end{pgfonlayer}

\end{tikzpicture}
\caption{Familias arquitectónicas derivadas del Transformer: \textit{encoder-only} para comprensión, \textit{decoder-only} para generación de texto, y \textit{encoder-decoder} para tareas de secuencia a secuencia. Elaboración propia basada en \citet{vaswani2017attention} y \citet{garcia2025desarrollo}.}
\label{fig:transformer_familias}
\end{figure}

El funcionamiento interno de los LLM se apoya en tres conceptos fundamentales que determinan su comportamiento \citep{garcia2025desarrollo}:

\begin{itemize}
  \item \textbf{Tokenización}: el texto se descompone en unidades mínimas (\textit{tokens}) mediante algoritmos como \textit{Byte Pair Encoding} (BPE), que manejan vocabularios de subpalabras y permiten representar cualquier idioma con un vocabulario finito y compacto.

  \item \textbf{Ventana de contexto}: es el número máximo de tokens que el modelo puede procesar en una sola solicitud. Esta capacidad ha crecido exponencialmente: de los 4\,096 tokens de GPT-3 (2020) a más de 1\,000\,000 en los modelos más recientes de 2025 --- GPT-5 supera el millón de tokens, Gemini~2.5~Pro llega a 2\,000\,000, y Llama~4~Scout alcanza los 10\,000\,000 de tokens de ventana nominal. Para el dominio de este trabajo, esto implica que modelos actuales pueden procesar la totalidad de una propuesta curricular en una sola consulta.

  \item \textbf{Leyes de escalado}: según el estudio Chinchilla \citep{garcia2025desarrollo}, existe una proporción óptima entre el número de parámetros del modelo y la cantidad de datos de entrenamiento. Modelos con 70\,000 millones de parámetros bien calibrados superan en rendimiento a modelos de mayor tamaño insuficientemente entrenados; el desempeño mejora de forma predecible al escalar parámetros, datos y cómputo de forma conjunta.
\end{itemize}

\subsection{Adaptación: entrenamiento, ajuste fino y técnicas de prompting}

El ciclo de vida de un LLM comprende tres etapas secuenciales \citep{torres2024desarrollo}, ilustradas en la Figura~\ref{fig:llm_ciclo_vida}:

\begin{figure}[H]
\centering
\begin{tikzpicture}[
  stbox/.style={rectangle, rounded corners=6pt,
                draw=#1!70, fill=#1!6,
                minimum width=3.6cm, minimum height=4.0cm,
                text width=3.4cm, align=center, inner sep=9pt},
  arr/.style={-{Stealth[length=5pt]}, very thick, color=gray!50}
]

%% Etapa 1 — Preentrenamiento
\node[stbox=Micolor1] (s1) at (0,0) {
  {\Large\bfseries\textcolor{Micolor1}{1}}\\[5pt]
  \textbf{Preentrenamiento}\\[6pt]
  \scriptsize Corpus masivos de texto general\\[3pt]
  \scriptsize (libros, webs, c\'{o}digo fuente)\\[8pt]
  \scriptsize\textit{Resultado:}\\
  \scriptsize capacidad ling\"{u}\'{\i}stica base
};

%% Etapa 2 — Ajuste fino
\node[stbox=teal] (s2) at (4.5,0) {
  {\Large\bfseries\textcolor{teal}{2}}\\[5pt]
  \textbf{Ajuste fino}\\
  \small\textit{(fine-tuning)}\\[6pt]
  \scriptsize Datos de dominio\\espec\'{\i}fico\\[3pt]
  \scriptsize Instrucci\'{o}n supervisada (SFT)\\[8pt]
  \scriptsize\textit{Resultado:}\\
  \scriptsize especializaci\'{o}n en tareas
};

%% Etapa 3 — Alineación RLHF
\node[stbox=gray] (s3) at (9.0,0) {
  {\Large\bfseries\textcolor{gray!70}{3}}\\[5pt]
  \textbf{Alineaci\'{o}n}\\
  \small\textit{(RLHF)}\\[6pt]
  \scriptsize Retroalimentaci\'{o}n humana\\[3pt]
  \scriptsize Reward model + PPO\\[8pt]
  \scriptsize\textit{Resultado:}\\
  \scriptsize respuestas \'{u}tiles y seguras
};

\draw[arr] (s1.east) -- (s2.west);
\draw[arr] (s2.east) -- (s3.west);

\end{tikzpicture}
\caption[Ciclo de vida de un LLM: tres etapas secuenciales.]{Ciclo de vida de un LLM. Elaboraci\'{o}n propia basada en \citet{torres2024desarrollo}.}
\label{fig:llm_ciclo_vida}
\end{figure}

En la fase de uso, la \textbf{ingenier\'{\i}a de prompts} (\textit{prompt engineering}) es la t\'{e}cnica que determina la calidad de las respuestas obtenidas. La Tabla~\ref{tab:llm_prompting} describe las estrategias m\'{a}s relevantes identificadas en la literatura.

\begin{table}[htb]
\centering
\small
\setlength{\extrarowheight}{4pt}
\caption{Principales técnicas de prompt engineering para LLM}
\label{tab:llm_prompting}
\begin{tabularx}{\textwidth}{p{2.4cm}XX}
\toprule
\textbf{Técnica} & \textbf{Descripción} & \textbf{Aplicación en este trabajo} \\
\midrule
Zero-shot         & El modelo responde sin ejemplos previos; solo con instrucciones en lenguaje natural                                & Generación de resúmenes y clasificación de secciones de propuestas de asignatura \\
\arrayrulecolor{gray!35}\hline
Few-shot          & Se incluyen 2--5 ejemplos en el prompt para guiar el formato y el estilo de la respuesta                          & Extracción de evidencias en formato CONEAU a partir de documentos institucionales \\
\hline
Chain of Thought (CoT) & El modelo descompone el razonamiento en pasos intermedios explícitos antes de responder                     & Evaluación de coherencia entre objetivos de aprendizaje y competencias de egreso \\
\hline
ReAct             & Combina razonamiento (\textit{Reasoning}) y acción (\textit{Acting}): el modelo intercala pasos de pensamiento con llamadas a herramientas externas & Orquestación de agentes que consultan la base documental, validan datos y generan reportes \\
\arrayrulecolor{black}\bottomrule
\end{tabularx}
\end{table}

\citet{garcia2025desarrollo} valida experimentalmente el patrón \textbf{ReAct} como la estrategia más efectiva para agentes que necesitan combinar razonamiento sobre el contexto con la invocación de herramientas externas. En este esquema, el LLM alterna entre generar \textit{trazas de pensamiento} y ejecutar acciones (llamadas a funciones, búsquedas, validaciones), refinando su respuesta tras recibir el resultado de cada acción.

\subsection{Generación Aumentada por Recuperación (RAG)}

Una limitación estructural de los LLM es que su conocimiento está congelado al momento del entrenamiento y no incluye documentación institucional privada. La técnica \textbf{RAG} (\textit{Retrieval-Augmented Generation}) resuelve este problema incorporando un pipeline de recuperación que alimenta al modelo con fragmentos de texto relevantes en el momento de la consulta, sin necesidad de re-entrenar el modelo \citep{farsane2025automatizacion}.

La Figura~\ref{fig:rag_pipeline} ilustra el pipeline RAG organizado en dos flujos complementarios: \textbf{indexación} (offline) y \textbf{recuperación} (online).

\begin{figure}[H]
\centering
\begin{tikzpicture}[
    box/.style={rectangle, rounded corners=4pt,
                draw=Micolor1!80, fill=white,
                minimum width=1.9cm, minimum height=0.95cm,
                align=center, font=\small\bfseries},
    sbox/.style={rectangle, rounded corners=4pt,
                draw=teal!70, fill=white,
                minimum width=1.9cm, minimum height=0.95cm,
                align=center, font=\small\bfseries},
    store/.style={cylinder, shape border rotate=90,
                draw=Micolor1!70, fill=white,
                minimum height=1.1cm, minimum width=1.4cm,
                align=center, font=\small\bfseries},
    tech/.style={font=\scriptsize, text=gray!60, align=center},
    arr/.style={-{Stealth[length=4.5pt]}, thick, color=Micolor1!80},
    sarr/.style={-{Stealth[length=4.5pt]}, thick, color=teal!70},
    varr/.style={-{Stealth[length=4.5pt]}, thick, color=gray!55, dashed},
    caplab/.style={font=\scriptsize\bfseries, align=center}]

  %--- Fila superior: Indexación ---
  \node[box]                         (doc)  at (0,    2.8) {Documentos};
  \node[box, right=0.45cm of doc]    (seg)  {Segmentaci\'{o}n};
  \node[box, right=0.45cm of seg]    (emb1) {Embeddings};
  \node[store, right=0.6cm of emb1]  (vs)   {Vector\\Store};
  \draw[arr] (doc)  -- (seg);
  \draw[arr] (seg)  -- (emb1);
  \draw[arr] (emb1) -- (vs);
  \node[tech, below=0.25cm of doc]  {\scriptsize PDF/DOCX};
  \node[tech, below=0.25cm of seg]  {\scriptsize Chunking};
  \node[tech, below=0.25cm of emb1] {\scriptsize text-embedding};

  %--- Fila inferior: Recuperación ---
  \node[sbox] (qry)  at (0,    0)   {Consulta};
  \node[sbox, right=0.45cm of qry]  (emb2) {Embedding};
  \node[sbox, right=0.45cm of emb2] (srch) {B\'{u}squeda\\sem\'{a}ntica};
  \node[box,  right=0.45cm of srch] (llm)  {LLM};
  \node[sbox, right=0.45cm of llm]  (resp) {Respuesta};
  \draw[sarr] (qry)  -- (emb2);
  \draw[sarr] (emb2) -- (srch);
  \draw[sarr] (srch) -- (llm);
  \draw[sarr] (llm)  -- (resp);
  \node[tech, below=0.25cm of qry]  {\scriptsize lenguaje\\natural};
  \node[tech, below=0.25cm of emb2] {\scriptsize mismo modelo};
  \node[tech, below=0.25cm of srch] {\scriptsize FAISS/Qdrant};
  \node[tech, below=0.25cm of llm]  {\scriptsize GPT\,/\,LLaMA};

  %--- Flecha vertical: Vector Store → Búsqueda ---
  \draw[varr] (vs.south) -- node[right,font=\scriptsize\itshape,text=gray!55]{top-k} (srch.north);

  %--- Fondos de grupo ---
  \begin{pgfonlayer}{background}
    \fill[Micolor1!7, rounded corners=6pt]
      ([xshift=-20pt, yshift=-18pt]doc.south west) rectangle
      ([xshift=18pt,  yshift=14pt]vs.north east);
    \fill[teal!7, rounded corners=6pt]
      ([xshift=-20pt, yshift=-18pt]qry.south west) rectangle
      ([xshift=18pt,  yshift=14pt]resp.north east);
  \end{pgfonlayer}

  %--- Etiquetas de capa ---
  \node[caplab, text=Micolor1!80]
    at ($(doc.north west) + (-0.1, 0.38cm)$) [anchor=west]
    {\textit{Flujo 1: Indexaci\'{o}n (offline)}};
  \node[caplab, text=teal!80]
    at ($(qry.north west) + (-0.1, 0.35cm)$) [anchor=west]
    {\textit{Flujo 2: Recuperaci\'{o}n (online)}};
\end{tikzpicture}
\caption{Pipeline RAG: flujo de indexación y flujo de recuperación. Elaboración propia basada en \citet{farsane2025automatizacion}.}
\label{fig:rag_pipeline}
\end{figure}

En el \textbf{flujo de indexación}, los documentos fuente se segmentan en fragmentos (\textit{chunks}), se codifican como vectores de alta dimensión mediante un modelo de embeddings, y se almacenan en una base vectorial (FAISS, Qdrant, ChromaDB). En el \textbf{flujo de recuperación}, la consulta del usuario se transforma con el mismo modelo de embeddings y se compara semánticamente contra el índice; los fragmentos más similares se recuperan y se inyectan en el prompt del LLM, que genera una respuesta fundamentada en el texto institucional \citep{farsane2025automatizacion}. Esta arquitectura provee \textbf{trazabilidad} (cada respuesta puede vincularse a su fuente documental) y permite actualizar el conocimiento sin re-entrenar el modelo.

\subsection{LLM en gestión documental especializada}

La literatura documenta una convergencia hacia el uso de LLM como núcleo semántico de sistemas de gestión documental en dominios regulados. Los trabajos más relevantes para este estudio abordan cuatro dimensiones complementarias: asistencia inteligente en lenguaje natural, automatización de workflows documentales, procesamiento de corpora normativos complejos y generación estructurada de documentos institucionales. La Tabla~\ref{tab:llm_doc_comparativa} sintetiza sus aportes principales.

\citet{garcia2025desarrollo} implementa un asistente basado en LLM que permite a usuarios no técnicos interactuar con conjuntos de datos mediante lenguaje natural, eliminando la necesidad de código. Sus resultados validan que el patrón \textbf{ReAct} supera en consistencia y cobertura a enfoques puramente conversacionales en dominios con terminología específica --- precisamente las condiciones del dominio CONEAU.

\citet{farsane2025automatizacion} demuestra que un sistema dual de workflows (indexación + recuperación semántica) mejora significativamente la precisión en la búsqueda de información interna y la coherencia entre fuentes heterogéneas. Su hallazgo clave: la \emph{segmentación semántica} de documentos --- dividirlos por unidades de significado y no por longitud fija --- incrementa la relevancia de los fragmentos recuperados, patrón que este trabajo incorpora en su pipeline RAG.

\citet{collado2024implementacion} aplica técnicas RAG sobre LLM a la extracción y generación automatizada de documentos en entidades públicas, demostrando que los modelos pueden producir salidas estructuradas conformes a plantillas institucionales cuando se los combina con un contexto documental preciso. El trabajo identifica la \emph{calidad del contexto inyectado} como el factor determinante de la precisión de las respuestas generadas.

\citet{torres2024desarrollo} valida el uso de modelos open-source (Mistral, LLaMA) para resumen, clasificación y búsqueda semántica sobre documentos legislativos de alta complejidad normativa en entornos con restricciones de privacidad. Sus conclusiones son transferibles al contexto universitario argentino: el despliegue local elimina dependencias de APIs externas y responde a las exigencias de soberanía de datos en organismos públicos regulados.

\citet{he2025llm} aporta una revisión sistemática de LLM aplicados a flujos de trabajo complejos, identificando como denominador común su capacidad para operar como \textbf{intermediarios semánticos} entre instrucciones en lenguaje natural y acciones sobre sistemas estructurados. Este atributo habilita que docentes, directores y secretarios académicos interactúen con la plataforma sin conocimiento técnico sobre las estructuras de datos subyacentes.

\begin{table}[H]
\centering
\small
\setlength{\extrarowheight}{3pt}
\caption{Trabajos representativos en LLM aplicados a gestión documental especializada}
\label{tab:llm_doc_comparativa}
\begin{tabularx}{\textwidth}{p{2.8cm} p{2.6cm} p{2.8cm} X}
\toprule
\textbf{Trabajo} & \textbf{Dominio} & \textbf{Técnica / Modelos} & \textbf{Aporte al presente trabajo} \\
\midrule
\citet{garcia2025desarrollo}
  & Análisis de datos en lenguaje natural
  & GPT + ReAct + herramientas
  & Valida ReAct para dominios con terminología específica \\
\arrayrulecolor{gray!35}\hline
\citet{farsane2025automatizacion}
  & Gestión del conocimiento empresarial
  & Pipeline RAG dual
  & Segmentación semántica como clave de precisión en recuperación \\
\hline
\citet{collado2024implementacion}
  & Gestión documental en entidades públicas
  & RAG + generación estructurada
  & Generación de documentos conformes a plantillas institucionales \\
\hline
\citet{torres2024desarrollo}
  & Documentos legislativos normativos
  & Mistral / LLaMA \mbox{open-source}
  & Modelos locales para instituciones con restricciones de privacidad \\
\hline
\citet{he2025llm}
  & Flujos de trabajo en ingeniería de software
  & LLM-MAS revisión sistemática
  & LLM como intermediarios semánticos entre lenguaje natural y sistemas \\
\arrayrulecolor{black}\bottomrule
\end{tabularx}
\end{table}

Estos trabajos establecen que los LLM son técnicamente maduros para asistir la gestión documental en dominios regulados: procesan vocabulario especializado, generan salidas estructuradas y escalan a grandes volúmenes de documentos. Ninguno, sin embargo, aborda documentación curricular bajo marcos de acreditación universitaria acreditadora.

\subsection{Brecha con el dominio CONEAU}

Los LLM han demostrado capacidad para procesar documentación regulada, extraer entidades especializadas, clasificar texto normativo y generar reportes estructurados. Sin embargo, \textbf{no se identifican en la literatura aplicaciones de LLM al procesamiento de propuestas de asignatura, planes de mejora ni evidencias de acreditación bajo el marco normativo de la CONEAU}. Los estándares de la Ordenanza 75/23 y los criterios de evaluación de pares exigen un vocabulario y una estructura semántica muy específicos que requieren adaptación deliberada del modelo. Este vacío refuerza la necesidad de integrar LLM con componentes de contexto institucional explícito ---como el RAG sobre documentos CONEAU--- elemento central del sistema propuesto en este trabajo.

\section{Model Context Protocol (MCP)}
\label{sec:eda_mcp}

El \textbf{Model Context Protocol (MCP)} es un protocolo abierto que estandariza la comunicación entre modelos de lenguaje de gran escala y sistemas externos --- herramientas, bases de datos, APIs y recursos documentales --- mediante una interfaz JSON-RPC orientada a sesión. Propuesto por Anthropic en noviembre de 2024, MCP resuelve la fragmentación que caracterizaba las integraciones previas de LLM con herramientas externas: cada integración requería código personalizado, controles de seguridad ad hoc y esquemas de contexto incompatibles entre plataformas \citep{ray2025survey, hou2025model, ehtesham2025survey}. En apenas doce meses, el protocolo se convirtió en el estándar de facto del ecosistema de agentes de IA, con más de 26\,000 servidores registrados públicamente y soporte oficial de OpenAI, Microsoft, Google y los principales entornos de desarrollo \citep{hou2025model}.

%%─────────────────────────────────────────────
\subsection{Origen y problemática que resuelve}
%%─────────────────────────────────────────────

Históricamente, la integración de LLM con sistemas externos evolucionó desde las interfaces simbólicas de la década de 1990 (KQML, FIPA-ACL) hasta los mecanismos de \textit{function calling} de OpenAI (2023) y los frameworks como LangChain y LlamaIndex. Sin embargo, todos estos enfoques compartían una limitación estructural: las definiciones de herramientas eran estáticas, los controles de acceso eran específicos de cada framework, y no existía un mecanismo de descubrimiento dinámico de capacidades \citep{ehtesham2025survey}.

\citet{ray2025survey} identifica tres barreras que MCP supera explícitamente: (1)~\textbf{interfaces sin estado} (\textit{stateless}), que impedían mantener contexto entre llamadas consecutivas; (2)~\textbf{controles de seguridad ad hoc}, no estandarizados ni auditables; y (3)~\textbf{ausencia de un mecanismo unificado} para inyectar contexto rico y multi-turno. El protocolo aborda estas limitaciones mediante un framework orientado a sesión, negociación de capacidades y control de acceso conforme a OAuth~2.1.

%%─────────────────────────────────────────────
\subsection{Arquitectura del protocolo}
%%─────────────────────────────────────────────

MCP define una arquitectura de tres capas: \textbf{Host}, \textbf{Cliente MCP} y \textbf{Servidor MCP}, cuyos roles se ilustran en la Figura~\ref{fig:mcp_arch} \citep{ehtesham2025survey, ray2025survey}.

\begin{figure}[H]
\centering
\begin{tikzpicture}[
  hostbox/.style={rectangle, rounded corners=7pt,
                  draw=Micolor1!80, fill=Micolor1!7,
                  minimum width=11.5cm, minimum height=1.5cm,
                  align=center, font=\small\bfseries},
  clientbox/.style={rectangle, rounded corners=6pt,
                    draw=teal!80, fill=teal!7,
                    minimum width=7.5cm, minimum height=1.3cm,
                    align=center, font=\small\bfseries},
  dbarr/.style={{Stealth[length=4pt]}-{Stealth[length=4pt]}, thick},
  lbl/.style={font=\scriptsize, text=gray!65, align=center}
]

%% --- HOST (top tier) ---
\node[hostbox] (host) at (0, 9.2) {
  \textcolor{Micolor1}{\normalsize LLM Application (Host)}\\[2pt]
  \small Claude Desktop\,·\,IDEs\,·\,Aplicaciones propias
};

%% --- CLIENT (middle tier) ---
\node[clientbox] (client) at (0, 7.0) {
  \textcolor{teal}{\normalsize MCP Client}\\[2pt]
  \small Enruta solicitudes\,·\,JSON-RPC\,2.0\,·\,Conexi\'{o}n 1:1 por servidor
};

%% --- SERVER A (bottom-left, built with tabular inside node) ---
\node[rectangle, rounded corners=5pt,
      draw=gray!60, fill=gray!5,
      minimum width=4.8cm, minimum height=4.0cm,
      align=center] (srv1) at (-3.2, 3.2)
  {\begin{tabular}{c}
     {\small\bfseries MCP Server A}\\[6pt]
     \colorbox{Micolor1!12}{\color{Micolor1}\scriptsize\bfseries\,Tools\,}\\[3pt]
     \colorbox{teal!12}{\color{teal}\scriptsize\bfseries\,Resources\,}\\[3pt]
     \colorbox{gray!20}{\scriptsize\bfseries\,Prompts\,}\\[3pt]
     \colorbox{orange!15}{\color{orange!80!black}\scriptsize\bfseries\,Sampling\,}
   \end{tabular}};

%% --- SERVER B (bottom-right) ---
\node[rectangle, rounded corners=5pt,
      draw=gray!60, fill=gray!5,
      minimum width=4.8cm, minimum height=4.0cm,
      align=center] (srv2) at (3.2, 3.2)
  {\begin{tabular}{c}
     {\small\bfseries MCP Server B}\\[6pt]
     \colorbox{Micolor1!12}{\color{Micolor1}\scriptsize\bfseries\,Tools\,}\\[3pt]
     \colorbox{teal!12}{\color{teal}\scriptsize\bfseries\,Resources\,}\\[3pt]
     \colorbox{gray!20}{\scriptsize\bfseries\,Prompts\,}\\[3pt]
     \colorbox{orange!15}{\color{orange!80!black}\scriptsize\bfseries\,Sampling\,}
   \end{tabular}};

%% --- ARROWS ---
\draw[dbarr, color=Micolor1!60, line width=1pt]
      (host.south) -- (client.north)
      node[midway, right=5pt, lbl] {instrucciones~/ respuestas};

\draw[dbarr, color=teal!60, line width=0.9pt]
      ($(client.south)+(-1.8, 0)$) -- (srv1.north)
      node[midway, left=3pt, lbl] {Stdio\\(local)};

\draw[dbarr, color=teal!60, line width=0.9pt]
      ($(client.south)+(1.8, 0)$) -- (srv2.north)
      node[midway, right=3pt, lbl] {HTTP+SSE\\(red)};

%% --- FOOTER NOTE ---
\node[lbl, text=gray!55] at (0, 0.8)
  {\scriptsize\textit{Cada capacidad (Tools, Resources, Prompts, Sampling) es accesible v\'{\i}a JSON-RPC\,2.0}};

\end{tikzpicture}
\caption[Arquitectura de tres capas del MCP \citep{ray2025survey, ehtesham2025survey}.]{Arquitectura de tres capas del Model Context Protocol: el \textit{Host} orquesta la aplicación LLM; el \textit{MCP Client} gestiona las conexiones con los servidores; cada \textit{MCP Server} expone cuatro capacidades (Tools, Resources, Prompts, Sampling) accesibles mediante JSON-RPC\,2.0 sobre Stdio (procesos locales) o HTTP+SSE (servidores remotos). Elaboración propia basada en \protect\citeauthor{ray2025survey} (\protect\citeyear{ray2025survey}) y \protect\citeauthor{ehtesham2025survey} (\protect\citeyear{ehtesham2025survey}).}
\label{fig:mcp_arch}
\end{figure}

El \textbf{Host} es la aplicación que aloja al LLM (p.~ej., Claude Desktop, un IDE con IA, o un agente personalizado). El \textbf{MCP Client}, embebido en el Host, establece y mantiene una conexión 1:1 con cada servidor, enviando solicitudes JSON-RPC~2.0 y recibiendo respuestas tipadas. Cada \textbf{MCP Server} es un proceso independiente --- local o remoto --- que expone sus capacidades al cliente. La capa de protocolo garantiza que cada solicitud esté vinculada a una respuesta predecible; la capa de transporte soporta comunicación local vía \textit{Stdio} y comunicación en red mediante HTTP con \textit{Server-Sent Events} (SSE) \citep{ehtesham2025survey}.

%%─────────────────────────────────────────────
\subsection{Capacidades del servidor MCP}
%%─────────────────────────────────────────────

Cada servidor MCP puede exponer cuatro tipos de capacidades, cada una con un modelo de control diferenciado (Tabla~\ref{tab:mcp_capacidades}):

\begin{table}[!h]
\centering
\small
\setlength{\extrarowheight}{3pt}
\caption{Capacidades del servidor MCP y su modelo de control}
\label{tab:mcp_capacidades}
\begin{tabular}{p{2.2cm} p{2.4cm} p{7.0cm} p{3.1cm}}
\toprule
\textbf{Capacidad} & \textbf{Control} & \textbf{Descripción} & \textbf{Ejemplo en este trabajo} \\
\midrule
\textbf{Tools}
  & Modelo (LLM)
  & El LLM invoca APIs o servicios externos, con o sin aprobación del usuario
  & Consultar y validar campos de una propuesta, ejecutar búsqueda semántica \\
\arrayrulecolor{gray!35}\hline
\textbf{Resources}
  & Aplicación
  & Documentos estructurados o datasets que el cliente selecciona y provee al LLM como contexto
  & Propuestas de asignatura, estándares de la Ordenanza 75/23, planes de mejora previos \\
\hline
\textbf{Prompts}
  & Usuario
  & Plantillas de interacción reutilizables que garantizan patrones de respuesta consistentes
  & Plantillas de evaluación por competencias y validación de propuestas \\
\hline
\textbf{Sampling}
  & Servidor
  & El servidor delega tareas de generación de texto al LLM del cliente
  & Generación automatizada de borradores de justificación curricular \\
\arrayrulecolor{black}\bottomrule
\end{tabular}
\end{table}

%%─────────────────────────────────────────────
\subsection{Ecosistema y adopción}
%%─────────────────────────────────────────────

Desde su lanzamiento en noviembre de 2024, el ecosistema MCP creció de manera exponencial. \citet{hou2025model} documentan, a septiembre de 2025, más de 26\,000 servidores en el directorio MCPWorld, 16\,592 en MCP.so y 13\,596 en MCP Servers Repository, entre otras plataformas. Los SDKs oficiales están disponibles en TypeScript, Python, Java, Kotlin y C\#. Las integraciones de mayor relevancia incluyen:

\begin{itemize}
  \item \textbf{OpenAI Agents SDK}: incorpora MCP para descubrimiento y llamada automática de herramientas, eliminando el código de conexión por herramienta \citep{hou2025model}.
  \item \textbf{Microsoft Copilot Studio}: anunció soporte oficial en marzo de 2025, habilitando integración unificada de herramientas empresariales.
  \item \textbf{Cursor / Windsurf / JetBrains}: los principales IDEs con asistencia de IA adoptaron MCP para conectar agentes de código con sistemas externos.
  \item \textbf{Claude Desktop} (Anthropic): implementación de referencia que popularizó el modelo Host-Client-Server.
\end{itemize}

%%─────────────────────────────────────────────
\subsection{Seguridad y trazabilidad}
%%─────────────────────────────────────────────

\citet{hou2025model} y \citet{ehtesham2025survey} identifican amenazas de seguridad específicas del ciclo de vida MCP que deben considerarse en implementaciones institucionales. Las más relevantes para el dominio de este trabajo son:

\begin{itemize}
  \item \textbf{Tool poisoning}: prompts o metadatos maliciosos que manipulan el comportamiento del LLM; mitigación mediante validación de esquemas y filtros semánticos.
  \item \textbf{Credential theft}: secretos filtrados a través de completaciones o salidas de herramientas; mitigación con OAuth~2.1 + PKCE y tokens de alcance restringido.
  \item \textbf{Command injection / RCE}: entradas no saneadas que disparan ejecución de sistema; mitigación deshabilitando acceso a shell y saneando entradas.
  \item \textbf{Tool redefinition} (\textit{rug pull}): herramientas que devienen maliciosas tras la validación inicial; mitigación con manifiestos firmados y versionados.
\end{itemize}

Para el presente trabajo, estas consideraciones se traducen en requisitos de diseño concretos: autenticación por roles institucionales, logs de auditoría inmutables y aislamiento de procesos entre agentes. La trazabilidad inherente al intercambio de mensajes JSON-RPC constituye, en este sentido, un activo técnico directo para demostrar cumplimiento normativo ante organismos auditores.

%%─────────────────────────────────────────────
\subsection{Posicionamiento comparativo}
%%─────────────────────────────────────────────

\citet{ehtesham2025survey} ubica al MCP en el contexto de cuatro protocolos emergentes de interoperabilidad para agentes de IA, proponiendo un mapa de adopción escalonada según el tipo de integración requerida (Tabla~\ref{tab:protocolos_agentes}):

\begin{table}[htb]
\centering
\small
\setlength{\extrarowheight}{3pt}
\caption{Comparación de protocolos de interoperabilidad para agentes de IA \citep{ehtesham2025survey}}
\label{tab:protocolos_agentes}
\begin{tabularx}{\textwidth}{p{1.8cm} p{1.6cm} p{2.0cm} X p{2.2cm}}
\toprule
\textbf{Protocolo} & \textbf{Origen} & \textbf{Transporte} & \textbf{Caso de uso principal} & \textbf{Modelo de seguridad} \\
\midrule
\textbf{MCP}
  & Anthropic (nov.\ 2024)
  & JSON-RPC / Stdio / HTTP+SSE
  & Acceso a herramientas y recursos desde un LLM
  & OAuth 2.1, tokens escopados \\
\arrayrulecolor{gray!35}\hline
\textbf{ACP}
  & IBM (2024)
  & RESTful HTTP, MIME multipart
  & Mensajería multimodal entre agentes heterogéneos
  & RBAC + DIDs \\
\hline
\textbf{A2A}
  & Google (abr.\ 2025)
  & HTTP + SSE, JSON-RPC
  & Delegación peer-to-peer entre agentes empresariales
  & Agent Cards firmadas, JWS \\
\hline
\textbf{ANP}
  & Comunidad (2024)
  & Protocolo abierto, JSON-LD
  & Colaboración descentralizada en internet abierto
  & W3C DIDs, cifrado E2E \\
\arrayrulecolor{black}\bottomrule
\end{tabularx}
\end{table}

La hoja de ruta propuesta por \citet{ehtesham2025survey} sugiere comenzar con MCP para acceso a herramientas, avanzar hacia ACP para mensajería estructurada multi-modal, incorporar A2A para ejecución colaborativa de tareas y finalmente extender a ANP para ecosistemas descentralizados. Para el sistema propuesto en este trabajo, MCP es la capa de partida natural: cubre el acceso a herramientas institucionales, la provisión de contexto documental y la generación de respuestas estructuradas, sin la complejidad adicional que A2A o ANP requieren.

%%─────────────────────────────────────────────
\subsection{Relevancia para el presente trabajo}
%%─────────────────────────────────────────────

La revisión realizada en las secciones anteriores permite establecer con claridad que los sistemas de gestión documental inteligente existentes --- incluso los más avanzados en términos de LLM y RAG --- adolecen de una limitación estructural: no proveen un mecanismo estandarizado para que los distintos componentes del sistema (agentes, bases documentales, módulos de validación, interfaces de usuario) se comuniquen, compartan contexto y sean auditados de manera uniforme \citep{ray2025survey}. En el dominio específico de la acreditación universitaria, donde cada decisión debe poder rastrearse hasta su fundamento normativo y donde actores con roles diferenciados interactúan sobre los mismos documentos, esta limitación no es menor: es arquitectónicamente bloqueante. El MCP es la respuesta técnica a ese problema.

La adopción del MCP como capa de interoperabilidad del sistema propuesto no es una elección de conveniencia, sino una decisión fundamentada en la correspondencia directa entre las propiedades del protocolo y los requisitos del dominio:

\begin{itemize}
  \item \textbf{Integración de componentes heterogéneos sin acoplamiento fuerte}: el sistema propuesto articula elementos tecnológicamente diversos --- una base documental con búsqueda semántica, agentes LLM especializados, módulos de validación normativa, un motor de flujo de trabajo y una interfaz web institucional ---. Sin MCP, cada par de componentes requeriría su propia integración ad hoc, con esquemas de datos incompatibles y lógica de conexión duplicada. MCP define una interfaz única, descubrible y versionada que permite reemplazar o escalar cualquier componente de forma independiente, exactamente como lo exige un sistema que debe evolucionar conforme cambia la normativa CONEAU \citep{ehtesham2025survey}.

  \item \textbf{Trazabilidad estructural como requisito de acreditación}: la Ordenanza~75/23 y los estándares CONEAU exigen que las instituciones puedan demostrar, con evidencia documental verificable, la coherencia entre el diseño curricular, las competencias declaradas, la bibliografía y las condiciones de dictado de cada asignatura \citep{ordenanza75}. Esto no es solo un requerimiento de reporte: es un requerimiento de trazabilidad de origen. Cada intercambio entre el Host y los servidores MCP genera mensajes JSON-RPC estructurados que registran qué herramienta fue invocada, con qué parámetros, en qué momento y por qué agente --- constituyendo un log de auditoría completo sin instrumentación adicional \citep{hou2025model}. Ningún framework de agentes previo garantiza esta propiedad de forma nativa.

  \item \textbf{Control de acceso granular alineado con la estructura institucional}: el dominio impone una jerarquía de roles con permisos diferenciados: docentes pueden crear y editar propuestas propias, directores de carrera pueden revisar y aprobar las de su departamento, secretarías académicas consolidan la información institucional, y pares evaluadores externos acceden en modo lectura durante el proceso de acreditación. El modelo de autorización OAuth~2.1 con tokens de alcance restringido nativo de MCP traduce esta jerarquía directamente en permisos sobre herramientas y recursos específicos, sin requerir una capa de autorización paralela \citep{ray2025survey}.

  \item \textbf{Extensibilidad hacia flujos multi-agente}: la arquitectura MCP habilita naturalmente la delegación de tareas entre agentes especializados --- un agente de validación normativa, un agente de consistencia bibliográfica, un agente de coherencia competencias-contenidos ---. Esta extensibilidad es clave para el diseño del sistema propuesto, que adopta una arquitectura de sistemas multi-agente (MAS) como paradigma central. MCP actúa aquí como el bus de comunicación que hace posible la coordinación inter-agente dentro de un mismo Host o entre servidores distribuidos \citep{ehtesham2025survey}.
\end{itemize}

\paragraph{Vacío identificado en la literatura.}
Una búsqueda sistemática en las principales bases de datos académicas (Scopus, Web of Science, IEEE Xplore, Google Scholar) no arroja ningún trabajo que aplique el Model Context Protocol al dominio de la educación superior, la gestión curricular ni los procesos de acreditación universitaria. Los estudios existentes sobre MCP se circunscriben a casos de uso en desarrollo de software, automatización de infraestructura en la nube, ciberseguridad y asistencia médica \citep{ray2025survey, hou2025model}. Tampoco se identifican propuestas que combinen MAS + LLM + RAG + MCP en ningún dominio institucional regulado.

Este vacío es significativo por dos razones. En primer lugar, confirma que la problemática de la acreditación universitaria no ha sido abordada con las herramientas tecnológicas contemporáneas más adecuadas. En segundo lugar, valida que la contribución técnica del presente trabajo --- diseñar e implementar un sistema que articule estas cuatro tecnologías en el contexto específico de la normativa CONEAU --- es genuinamente novedosa y no replicable desde trabajos previos. La originalidad del enfoque no reside en ninguna de las tecnologías individualmente, sino en su combinación deliberada para resolver un problema concreto, complejo y socialmente relevante: asistir a las instituciones universitarias argentinas en la construcción de evidencias de calidad auditables para sus procesos de acreditación.

\section{Síntesis y posicionamiento}
\label{sec:eda_sintesis}

La revisión realizada en este capítulo recorre cuatro dominios de conocimiento que convergen en la propuesta de este trabajo: los marcos normativos de la acreditación universitaria argentina, los sistemas multi-agente como paradigma de coordinación distribuida, los modelos de lenguaje de gran escala como capa de inteligencia semántica, y el Model Context Protocol como estándar de interoperabilidad y trazabilidad. El siguiente diagnóstico sintetiza los hallazgos:

\begin{itemize}
    \item \textbf{Dominio regulatorio}: los procesos de acreditación argentina son complejos, iterativos y altamente dependientes de documentación institucional sin soporte tecnológico específico. La Ordenanza 75/23 y los estándares de la CONEAU establecen requisitos precisos sobre estructura, coherencia y evidencia de las propuestas docentes, pero su verificación continúa siendo manual y propensa a omisiones \citep{ordenanza62, ordenanza75, duarte2023literature, sanchez2022acreditacion}.

    \item \textbf{Gestión documental inteligente}: los avances recientes en extracción de información, RAG y sistemas de gestión basados en LLM demuestran que la automatización es técnicamente viable para dominios documentales complejos \citep{cascavita2025sistema, paun2025implementacion, cruz2024diseno}. Sin embargo, ninguna de estas soluciones está diseñada para el marco normativo de la CONEAU ni contempla el ciclo de vida completo de una propuesta institucional regulada.

    \item \textbf{Sistemas multi-agente}: la arquitectura MAS provee los mecanismos de coordinación, delegación y especialización necesarios para orquestar los múltiples actores y procesos del ciclo de acreditación. Los trabajos revisados confirman su adecuación para dominios con roles diferenciados y flujos de trabajo complejos \citep{perez2022breve, sakurada2021multi}, aunque ninguno los aplica al gobierno de propuestas curriculares.

    \item \textbf{Modelos de lenguaje}: los LLM habilitan la asistencia semántica sobre documentación regulada, la detección de inconsistencias en lenguaje natural y la generación de sugerencias contextualmente adecuadas para usuarios docentes no técnicos \citep{garcia2025desarrollo, farsane2025automatizacion, agashe2023llm}. Su integración en flujos editoriales institucionales es, sin embargo, incipiente.

    \item \textbf{Model Context Protocol}: el MCP provee la capa de interoperabilidad, descubrimiento de capacidades y trazabilidad estructural que distingue a un sistema institucional auditado de una integración ad hoc. Su ecosistema crece aceleradamente, pero ningún trabajo lo aplica al dominio de la educación superior \citep{ehtesham2025survey, hou2025model, ray2025survey}.
\end{itemize}


El estado del arte evidencia que cada uno de estos componentes ha sido estudiado y validado de forma independiente. Lo que \textbf{no existe} es una propuesta que los articule como sistema integrado orientado a la acreditación universitaria regulada. En particular, la literatura no ofrece soluciones que combinen: gestión del ciclo de vida completo de propuestas de asignatura (desde la elaboración hasta la aprobación institucional), control editorial diferenciado por roles (docente, director de carrera, secretaría de carrera, miembro de la comisión curricular), asistencia inteligente alineada con los requisitos normativos de la CONEAU, y trazabilidad auditable de evidencias mediante un protocolo estandarizado.

Esta convergencia de ausencias no es casual. Refleja que el problema de la acreditación universitaria argentina no ha sido abordado desde la perspectiva de la ingeniería de sistemas inteligentes, y que las soluciones disponibles en dominios adyacentes no son directamente transferibles sin una reingeniería profunda orientada al marco normativo específico.

Este trabajo se posiciona en la intersección de esos cuatro dominios para cubrir la brecha identificada. La combinación MAS + LLM + RAG + MCP, aplicada al ciclo de vida de propuestas docentes bajo los estándares de la Ordenanza 75/23 y la CONEAU, constituye una contribución técnica y metodológica que ningún trabajo previo ofrece. 

